\setcounter{section}{-1}
\section{Preliminaries}\label{sec:prelim}

\subsection{Notation}
Graphs are finite and simple, $|G|=|V(G)|$, $N(v)$ is the neighbourhood of $v$, and $G[X]$ is the
subgraph induced on $X\subseteq V(G)$. We freely identify $X$ with $G[X]$. For disjoint $A,B$, $e(A,B)$ is
the number of edges between them. $A$ is \emph{complete} (\emph{anticomplete}) to $B$ if every (no) pair
$ab$ with $a\in A$, $b\in B$ is an edge. The pair is \emph{pure} if it is one of the two. A vertex is
\emph{mixed} on a set if it has both a neighbour and a non-neighbour there.

\begin{definition}[Sparsity]
Let $x\ge0$. A graph $G$ is \emph{$x$-sparse} if every vertex has at most $x|G|$ neighbours, and
\emph{$x$-restricted} if $G$ or $\Gbar$ is $x$-sparse. A set $X$ is $x$-sparse or $x$-restricted if $G[X]$
is. For disjoint $A,B$: $B$ is \emph{$x$-sparse to} $A$ if every vertex of $B$ has at most $x|A|$
neighbours in $A$, and $(A,B)$ is \emph{weakly $x$-sparse} if $e(A,B)\le x|A||B|$. If $B$ is $x$-sparse to
$A$ then $(A,B)$ is weakly $x$-sparse.
\end{definition}

\begin{definition}[Blockades]
For reals $k,w$, a \emph{$(k,w)$-blockade} is a sequence $(B_1,\dots,B_m)$ of pairwise disjoint sets
(\emph{blocks}, possibly empty) with $m\ge k$ and $|B_i|\ge w$ for all $i$. It is \emph{complete},
\emph{anticomplete} or \emph{pure} if every two blocks are complete, anticomplete or pure, respectively.
It is \emph{$x$-sparse} if $B_j$ is $x$-sparse to $B_i$ whenever $i<j$, and \emph{$x$-semisparse} if every
two blocks are complete or weakly $x$-sparse. A pure blockade is $0$-semisparse; an $x$-sparse blockade is
$x$-semisparse.
\end{definition}

A set $X$ is \emph{anticonnected} if $\Gbar[X]$ is connected, equivalently if every partition of $X$ into
two nonempty parts has a non-adjacent pair across. An \emph{anticomponent} of $X$ is the vertex set of a
component of $\Gbar[X]$. Distinct anticomponents of $X$ are complete to each other, and each anticomponent
is complete to the rest of $X$. We write $\hom(G)=\max(\alpha(G),\omega(G))=\hom(\Gbar)$. $P_k$ is the
path on $k$ vertices, $\cP$ is its complement for $k=6$, namely six vertices whose five non-adjacent pairs
form a path, and the \emph{house} is $\overline{P_5}$. Logarithms are to base $2$.

\subsection{The four imported theorems}\label{sec:imported}
We quote each result as stated in its source, together with the definitions it depends on, and then give
the Lean proposition that renders it (file \texttt{formal/EHP6/Imported.lean}). The main formal theorem
takes these four propositions as hypotheses. Theorem~\ref{imp:path} is proved in the formalization, so
the unconditional formal result assumes only the other three; see Section~\ref{sec:formal}. Each proposition is stated for
an arbitrary vertex set $S$ of a graph $G$, meaning the induced subgraph $G[S]$. This is equivalent to the source statement because induced
subgraphs of $H$-free graphs are $H$-free. In the Lean code, \lean{Free G H} means that $G$ has no induced
copy of $H$, \lean{Sparse G x S} means that every vertex of $S$ has at most $x|S|$ neighbours in $S$, and
\lean{Restricted G x S} means \lean{Sparse G x S} or \lean{Sparse Gᶜ x S}. Section~\ref{sec:formal} gives the
full definitions.

\begin{importedthm}[R\"odl~\cite{Rodl}; stated as {\cite[Theorem 1.3]{NSS7}}]\label{imp:rodl}
\emph{Source:} ``For every $\eps\in(0,1/2)$ and every graph $H$, there exists $\delta>0$ such that every
$H$-free graph $G$ has an $\eps$-restricted induced subgraph with at least $\delta|G|$ vertices.'' Here
$G$ is \emph{$\eps$-sparse} if its maximum degree is at most $\eps|G|$, and \emph{$\eps$-restricted} if one
of $G,\Gbar$ is $\eps$-sparse. We use it only for $H=\cP$ and write $\delta(\eps)$ for the constant.
\begin{lstlisting}
def RodlCoP6 : Prop := ∀ ε : ℝ, 0 < ε → ε < 1 / 2 → ∃ δ : ℝ, 0 < δ ∧
  ∀ (V : Type) [Fintype V] [DecidableEq V] (G : SimpleGraph V) [DecidableRel G.Adj],
    Free G P6ᶜ → ∀ S : Finset V, ∃ F ⊆ S, δ * S.card ≤ F.card ∧ Restricted G ε F
\end{lstlisting}
\end{importedthm}

\begin{importedthm}[{\cite[statement 3.1]{NSS5}}]\label{imp:path}
\emph{Source:} ``Let $P$ be a path with $k\ge2$ vertices, and let $0<y\le1/(60k)$. Let $G$ be a
$y^2$-sparse graph. Then either $G$ contains a copy of $P$, or $G$ contains an anticomplete
$(1/y,\lfloor y^2|G|\rfloor)$-blockade in $G$.'' In~\cite{NSS5} a blockade is a finite sequence of possibly
empty disjoint sets, and a $(k,w)$-blockade has length at least $k$ and width at least $w$, as above. A
``copy'' of $P$ is an induced copy: $H$-freeness in~\cite{NSS5} refers to induced subgraphs, and the
proof of~\cite[3.1]{NSS5} constructs an induced path. We use it for $P=P_6$, so $k=6$ and $y\le1/360$.
\begin{lstlisting}
def NssPath6 : Prop := ∀ y : ℝ, 0 < y → y ≤ 1 / 360 →
  ∀ (V : Type) [Fintype V] [DecidableEq V] (G : SimpleGraph V) [DecidableRel G.Adj]
    (S : Finset V), Sparse G (y ^ 2) S → ¬ ContainsInduced G P6 S →
      ∃ β : Blockade S (1 / y) (⌊y ^ 2 * S.card⌋₊ : ℝ), β.IsAnticomplete G
\end{lstlisting}
\emph{Proved in the formalization} (\texttt{formal/EHP6/NSSPath.lean}, for every $k\ge1$), following the
proof in~\cite{NSS5}. The only change is that the repeated removal of small sets with few neighbours
is replaced by a maximal family of such sets, which are pairwise disjoint and anticomplete; the maximality
gives the expansion property of the remaining vertices directly.
\end{importedthm}

\begin{importedthm}[Nguyen, Scott and Seymour~{\cite[Theorem 1.2]{NSS7}}]\label{imp:eh5}
\emph{Source:} ``$P_5$ satisfies the Erd\H{o}s--Hajnal conjecture,'' that is, there exists $c>0$ such that
every $P_5$-free graph $G$ has a clique or a stable set of size at least $|G|^{c}$. We fix such a
$\kappa>0$. Since $\hom(G)=\hom(\Gbar)$, every house-free graph $G$ also has $\hom(G)\ge|G|^{\kappa}$.
\begin{lstlisting}
def EhP5 : Prop := ∃ c : ℝ, 0 < c ∧
  ∀ (V : Type) [Fintype V] [DecidableEq V] (G : SimpleGraph V) [DecidableRel G.Adj],
    Free G P5 → ∃ T : Finset V, (IsCliqueF G T ∨ IsStableF G T) ∧ (Fintype.card V : ℝ) ^ c ≤ T.card
\end{lstlisting}
\end{importedthm}

\begin{importedthm}[the comb lemma; {\cite[Lemma 4.3]{NSS7}}, a special case of a lemma of~\cite{C5}]\label{imp:comb}
\emph{Source:} ``Let $G$ be a graph and let $A,B\subseteq V(G)$ be nonempty and disjoint, such that each
vertex in $A$ has at most $\Delta>0$ neighbours in $B$. Then either at most $20\sqrt{|B|\Delta}$ vertices
in $B$ have a neighbour in $A$; or for some integer $k\ge1$, there is a $(k,|B|/k^2)$-comb
$((a_i,B_i):i\in[k])$ in $G$ where $a_i\in A$ and $B_i\subset B$ for all $i\in[k]$.'' Here an
\emph{$(\ell,w)$-comb} is a sequence of pairs $((a_i,B_i):i\in[\ell])$ where $(B_1,\dots,B_\ell)$ is an
$(\ell,w)$-blockade, the $a_i$ are pairwise distinct, $\{a_1,\dots,a_\ell\},B_1,\dots,B_\ell$ are pairwise
disjoint, and for distinct $i,j$, $a_i$ is complete to $B_i$ and anticomplete to $B_j$. The symbol
$\subset$ in the source denotes inclusion; for $k=1$ it must allow $B_1=B$.
\begin{lstlisting}
def NssComb : Prop :=
  ∀ (V : Type) [Fintype V] [DecidableEq V] (G : SimpleGraph V) [DecidableRel G.Adj]
    (A B : Finset V) (Δ : ℝ), Disjoint A B → A.Nonempty → B.Nonempty → 0 < Δ →
    (∀ a ∈ A, ((nbrs G a B).card : ℝ) ≤ Δ) →
    (((B.filter (fun b => ∃ a ∈ A, G.Adj a b)).card : ℝ) ≤ 20 * Real.sqrt (B.card * Δ)) ∨
    (∃ ℓ : ℕ, 1 ≤ ℓ ∧ ∃ (a : Fin ℓ → V) (C : Fin ℓ → Finset V),
      Function.Injective a ∧ (∀ i, a i ∈ A) ∧ (∀ i, C i ⊆ B) ∧
      (∀ i j, i ≠ j → Disjoint (C i) (C j)) ∧
      (∀ i, (B.card : ℝ) / ℓ ^ 2 ≤ (C i).card) ∧
      (∀ i, ∀ c ∈ C i, G.Adj (a i) c) ∧
      (∀ i j, i ≠ j → ∀ c ∈ C j, ¬ G.Adj (a i) c))
\end{lstlisting}
The disjointness of the apexes from the teeth follows from $A\cap B=\emptyset$.
\end{importedthm}

\subsection{Two elementary lemmas of Nguyen, Scott and Seymour}
The proof also uses Lemmas 4.1 and 4.2 of~\cite{NSS7}. We prove both, since they are short and are
proved in the formalization. Lemma~\ref{lem:locate} is~\cite[Lemma 4.2]{NSS7} in the form of arXiv version~3
(February 2026), with the added hypothesis $A\ne\emptyset$: for $A=\emptyset\ne B$ the hypothesis holds
vacuously but the conclusion fails. This does not affect any use of it, here or in~\cite{NSS7}. The
revision of~\cite{NSS7} dated July 2026 states the lemma for all $x>0$ with the weaker bound
$|A'|\le2/x$; since we prove Lemma~\ref{lem:locate} below, nothing here depends on which version is used.

\begin{lemma}[{\cite[Lemma 4.1]{NSS7}}]\label{lem:cores}
Let $k\ge2$ be an integer and let $G$ be a graph all of whose anticomponents have fewer than $|G|/k$
vertices. Then $G$ has a complete $(k,|G|/k^2)$-blockade.
\end{lemma}
\begin{proof}
Let $n=|G|$, $w=n/k^2$, $c=n/k$. Distinct anticomponents are complete to each other, so if
$\cG_1,\dots,\cG_k$ are disjoint collections of anticomponents, the unions $\bigcup\cG_1,\dots,\bigcup\cG_k$
are pairwise complete. From any
collection of anticomponents of total size at least $w$ we can pick a subcollection of total size in
$[w,w+c)$: add members one at a time until the total reaches $w$, and use that every member has size
$<c$. Pick such subcollections $\cG_1,\cG_2,\dots$ greedily from the anticomponents not yet used. After
$j\le k-1$ steps at most $(k-1)(w+c)=n-n/k^2$ vertices are used, so the remaining anticomponents have
total size at least $w$ and step $j+1$ is possible. The unions of $\cG_1,\dots,\cG_k$ form a complete
$(k,n/k^2)$-blockade.
\end{proof}

Applied to $\Gbar$: if every component of $G$ has fewer than $|G|/k$ vertices, then $G$ has an
anticomplete $(k,|G|/k^2)$-blockade.

\begin{lemma}[{cf.~\cite[Lemma 4.2]{NSS7}}]\label{lem:locate}
Let $x\in(0,\frac12)$, let $A\ne\emptyset$ and $B$ be sets of vertices, and suppose every vertex of $B$ has
at least $x|A|$ neighbours in $A$. Then there is $A'\subseteq A$ with $|A'|\le1/x$ such that at least
$|B|/2$ vertices of $B$ have a neighbour in $A'$.
\end{lemma}
\begin{proof}
Let $U\subseteq B$. Counting edges, $\sum_{a\in A}|N(a)\cap U|=\sum_{b\in U}|N(b)\cap A|\ge x|A||U|$, so
some $a\in A$ has at least $x|U|$ neighbours in $U$. Choose $a_1,a_2,\dots,a_t$ greedily, each covering an
$x$-fraction of the vertices of $B$ not yet covered, where $t=\lfloor1/x\rfloor$. At most $(1-x)^t|B|$
vertices remain uncovered. As $x<\frac12$, $t\ge2$ and $x>1/(t+1)$, so $1-x<t/(t+1)$, and Bernoulli's
inequality $(1+1/t)^t\ge2$ gives $(1-x)^t\le(t/(t+1))^t\le\frac12$. Take $A'=\{a_1,\dots,a_t\}$.
\end{proof}

\begin{lemma}\label{lem:anticonn}
If $D$ is anticonnected and $1\le b\le|D|$, then $D$ contains an anticonnected set of exactly $b$
vertices.
\end{lemma}
\begin{proof}
Start from a single vertex. If $B\subsetneq D$ is anticonnected, anticonnectivity of $D$ applied to the
partition $(B,D\setminus B)$ gives $p\in B$ and $q\in D\setminus B$ with $pq\notin E(G)$, and $B\cup\{q\}$
is anticonnected.
\end{proof}

\subsection{Modular decomposition}\label{sec:md}
Let $Y$ be a set of vertices. A \emph{module} of $G[Y]$ is a set $M\subseteq Y$ such that every vertex of
$Y\setminus M$ is complete or anticomplete to $M$. A module is \emph{strong} if it overlaps no other
module, where two sets overlap if they meet and neither contains the other. The following facts are
standard (Gallai~\cite{Gallai}); the forms used below are proved in the formalization. The strong modules
of $G[Y]$ form a laminar family containing $Y$ and all singletons, that is, a rooted tree under inclusion.
The \emph{children} of a strong module $N$ with $|N|\ge2$ are the maximal strong modules properly
contained in $N$, and they partition $N$. There are three kinds of node:
\begin{itemize}[leftmargin=1.5em]
\item $N$ is \emph{parallel} if $G[N]$ is disconnected; then its children are the components of $G[N]$;
\item $N$ is \emph{series} if $\Gbar[N]$ is disconnected; then its children are the anticomponents of
  $N$;
\item $N$ is \emph{prime} if $G[N]$ and $\Gbar[N]$ are both connected; then the quotient graph on the
  children (two children adjacent if complete to each other) is prime, that is, it has at least four
  vertices and no modules other than the trivial ones.
\end{itemize}
A module of $G[N]$ for a module $N$ of $G[Y]$ is a module of $G[Y]$.

\begin{lemma}[MD1]\label{lem:md1}
Disjoint modules $A,B$ of $G[Y]$ are pure to each other.
\end{lemma}
\begin{proof}
Each $a\in A$ lies outside $B$, so it is complete or anticomplete to $B$. If $a_1\in A$ were complete and
$a_2\in A$ anticomplete to $B$, any $b\in B$ would be mixed on the module $A$, impossible as $b\notin A$.
\end{proof}

\begin{lemma}[MD2]\label{lem:md2}
If $N$ is a prime node, every module $X$ of $G[N]$ equals $N$ or is contained in one child of $N$.
\end{lemma}
\begin{proof}
Let $I$ be the set of children meeting $X$, and suppose $|I|\ge2$. For a child $C\notin I$ and $c\in C$,
the adjacency of $c$ to a child $M\in I$ equals its adjacency to $X\cap M$ ($M$ is a module and
$c\notin M$), which equals its adjacency to $X$ ($X$ is a module and $c\notin X$). Hence $I$ is a module
of the prime quotient, and since $|I|\ge2$ it contains all children. If some child $M_1\not\subseteq X$,
take $m\in M_1\setminus X$. For every other child $M$ the adjacency of $m$ to $M$ equals its adjacency to
$X\cap M$, which equals its adjacency to $X$. So in the quotient, $M_1$ is complete or anticomplete to all
other children, and these, at least three of them, form a nontrivial module of the prime quotient, a
contradiction. Hence every child lies in $X$, and $X=N$.
\end{proof}
